En los \'ultimos a\~nos, la electromovilidad se ha consolidado como uno de los pilares de la transformaci\'on hacia sistemas de transporte p\'ublico m\'as limpios y sostenibles, en un contexto de creciente preocupaci\'on por el cambio clim\'atico y la calidad del aire en las ciudades. Chile no ha sido ajeno a este proceso. Desde la llegada de los primeros buses el\'ectricos a Santiago a fines de 2018 \cite{worldbank2020}, el pa\'is ha avanzado hacia la electrificaci\'on de su transporte p\'ublico para enfrentar el hist\'orico problema de contaminaci\'on atmosf\'erica de la capital, reduciendo emisiones y ruido. Sin embargo, la incorporaci\'on de flotas el\'ectricas no es un reemplazo directo de la tecnolog\'ia convencional, ya que introduce restricciones operacionales nuevas que exigen repensar c\'omo se planifica la operaci\'on diaria. La Red Metropolitana de Movilidad (RED) de Santiago es uno de los sistemas de transporte p\'ublico m\'as grandes de Am\'erica Latina y cuenta con la flota de buses el\'ectricos m\'as grande fuera de China, lo que convierte a Santiago en un caso de estudio a escala real, con la complejidad operacional propia de una red metropolitana masiva.

La complejidad de incorporar buses el\'ectricos a un sistema planificado para buses di\'esel radica en que estos veh\'iculos no operan bajo la misma l\'ogica que sus predecesores. Un bus el\'ectrico tiene una autonom\'ia considerablemente menor por la limitaci\'on de su bater\'ia, que necesita tiempos de recarga considerables (de noche en los dep\'ositos o en pausas operacionales durante el d\'ia) y solo puede recargarse en electroterminales, instalaciones a\'un escasas y con pocos cargadores. Esto convierte la planificaci\'on en un problema conjunto de transporte y energ\'ia, donde adem\'as de asignar buses a los viajes debe decidirse cu\'ando y d\'onde recargar cada veh\'iculo sin saturar una infraestructura compartida y limitada. Frente a esta complejidad, la Investigaci\'on Operativa ofrece las herramientas adecuadas para abordar el problema mediante la optimizaci\'on matem\'atica, las heur\'isticas, la simulaci\'on y el an\'alisis de datos que permiten explorar sistem\'aticamente distintas estrategias de operaci\'on y evaluar su costo y desempe\~no.

En el resto de este informe se presenta una primera aproximaci\'on al problema de planificaci\'on operacional de la flota el\'ectrica de RED, donde se describe en detalle el problema abordado, se analizan los datos disponibles junto con los indicadores de desempe\~no considerados y se discuten las metodolog\'ias de resoluci\'on candidatas.
